\documentclass[../../main.tex]{subfiles}
\begin{document}

\begin{flushright}
\footnotesize\textit{【自動文字起こし・要確認】}
\end{flushright}

{\bf [解]}
題意より
\begin{align}
0<a<1, \qquad 0<c\le b \label{eq:1}
\end{align}
である．題意の関数を
\begin{align}
 f(x)=x^3+ax^2+bx+c \label{eq:2}
\end{align}
とすると$f(x)=0$の3実解が$x=\alpha,\beta,\gamma$である．対称性から
\begin{align}
 \alpha < \beta < \gamma \label{eq:3}
\end{align}
として良い．
因数定理から
\begin{align}
 f(x)=(x-\alpha)(x-\beta)(x-\gamma) \label{eq:4}
\end{align}
だから，\cref{eq:2,eq:3}の係数を比較して
\begin{align}
a=-(\alpha+\beta+\gamma), \qquad b=\beta\gamma+\gamma\alpha+\alpha\beta, \qquad c=-\alpha\beta\gamma \label{eq:5}
\end{align}
を得る．これを\cref{eq:1}に代入して
\begin{align}
-1&<\alpha+\beta+\gamma<0  \label{eq:6}\\
\alpha\beta\gamma&<0       \label{eq:7}\\
0&\le\beta\gamma+\gamma\alpha+\alpha\beta+\alpha\beta\gamma \label{eq:8}
\end{align}
を得る．

\cref{eq:7}より$\alpha,\beta,\gamma$のうちの奇数個が負であって，これと$\alpha<\beta<\gamma$より
\begin{itemize}
 \item[(i)] $\alpha<0<\beta<\gamma$ \\
 \item[(ii)] $\alpha<\beta<\gamma<0$
\end{itemize}
のいずれかである．以下これに従って場合分けする．

\subparagraph{(i) $\alpha<0<\beta<\gamma$の時}

$\alpha,\beta,\gamma$がすべて負で$\alpha$が最小だから\cref{eq:6}において
\begin{align}
 -1<\alpha+\beta+\gamma < \alpha \quad \therefore -1 < \alpha
\end{align}
だから
\begin{align}
 -1 < \alpha < \beta < \gamma < 0
\end{align}
であり題意を満たす．


\subparagraph{(ii) $\alpha<0<\beta<\gamma$の時}

$p=\beta+\gamma,\ q=\beta\gamma$とおくと，$\beta,\gamma$は$t$の2次式$t^2-pt+q=0$の正の2実解だから，
\begin{align}
p^2-4q>0, \qquad p>0,\ q>0 \label{eq:9}
\end{align}
である．$p,q$を\cref{eq:6,eq:8}に代入して
\begin{align}
-1-\alpha<p<-\alpha, \qquad 0\le\alpha p+(1+\alpha)q \label{eq:10}
\end{align}
である．まず\cref{eq:10}の第二式が$\alpha<0, p,q>0$のもとで成り立つには
\begin{align}
 1+\alpha > 0 \label{eq:11}
\end{align}
であって，このもとで\cref{eq:9,10}は
\begin{align}
 \begin{cases}
  0<p<-\alpha,\ q>0 \\
 \dfrac{-\alpha}{1+\alpha}p \le q < \dfrac{p^2}{4}
 \end{cases} \label{eq:12}
\end{align}
となる．しかし$q=\dfrac{-\alpha}{1+\alpha}p$と$q=\dfrac{p^2}{4}$の$p>0$での交点である$p=\dfrac{-4\alpha}{1+\alpha}$は$p=-\alpha$より大きいため，下図のようにこれを満たすような$(p,q)$は存在せず，よってこの場合は不適である．


\begin{figure}[htb]
\begin{center}
\begin{tikzpicture}[scale=1.2, >=stealth]

  % --- 1. 定数・座標の設定 ---
  % alpha = -0.5 の例 (1+alpha = 0.5)
  % 傾き k = -alpha/(1+alpha) = 0.5/0.5 = 1
  % 交点 p_int = -4*alpha/(1+alpha) = 4
  % 境界 p_bound = -alpha = 0.5
  
  \draw[->, thick] (-0.5,0) -- (5.2,0) node[right] {$p$};
  \draw[->, thick] (0,-0.5) -- (0,4.8) node[above] {$q$};
  \node[below left] at (0,0) {$O$};

  % --- 2. 領域条件: 0 < p < -alpha (不適合を示す対象の帯状領域) ---
  % p = -alpha の境界（x=0.5）
  \fill[red!15, opacity=0.7] (0,0) rectangle (1.0,4.2);
  \draw[dashed, red!80!black, thick] (1.0,0) -- (1.0,4.2) 
    node[above, font=\small] {$p = -\alpha$};

  % --- 3. 放物線 q = p^2 / 4 の描画 ---
  % 放物線の下側 (q < p^2/4) が条件の1つ
  \draw[domain=0:4.2, samples=100, ultra thick, blue!80!black] 
    plot (\x, {0.25*\x*\x}) node[right] {$q = \frac{p^2}{4}$};

  % --- 4. 直線 q = {-alpha / (1+alpha)} * p の描画 ---
  % 直線の上側 (q >= ... p) がもう1つの条件
  \draw[domain=0:4.2, samples=2, ultra thick, green!50!black] 
    plot (\x, {0.9*\x}) node[above right] {$q = \frac{-\alpha}{1+\alpha}p$};

  % --- 5. 2曲線の交点 (p = -4*alpha / (1+alpha)) ---
  % 傾き 0.9 のとき交点は p = 3.6
  \coordinate (Pint) at (3.6, {0.25*3.6*3.6});
  \fill[black] (Pint) circle (2pt);
  \draw[dotted, gray] (3.6,0) -- (Pint) -- (0, {0.25*3.6*3.6});
  \node[below] at (3.6,0) {$\frac{-4\alpha}{1+\alpha}$};

  % --- 6. 条件を満たす領域の存在不可能性を示す強調 ---
  % 0 < p < -alpha の領域内で「直線の上側」かつ「放物線の下側」となる領域がないことを矢印等で明示
  \draw[<->, thick, red!80!black] (0.2, 0.18) -- (0.2, 0.01) 
    node[midway, right, font=\tiny] {空集合};

  % 帯状領域の注釈
  \node[red!80!black, font=\small, align=center] at (0.5, 3.2) 
    {$0 < p < -\alpha$\\[1pt]の範囲};

\end{tikzpicture}
\caption{\cref{eq:12}の領域}
\end{center} 
\end{figure}


\casesend

以上(i),(ii)の場合分けから題意は示された．

\newpage

\end{document}
